% !TEX root = ../main.tex
\section{Conclusion}

This study reproduced and extended the central result of \citet{daugherty1991credibility}: open-loop LP forest plans with inter-period harvest-flow constraints are frequently dynamically inconsistent under sequential replanning: over the thesis's observation window, 49\% of our flow-constrained scenarios exceed a 5\% divergence tolerance, although on the thesis's own volume measure the magnitudes are smaller than the thesis reported (5.3\% over the factor combinations we could run vs 9.0\% over all of the thesis's runs; \citealp[Table~6.2]{daugherty1991credibility}), and the thesis's patterns by policy and by landbase are only partly reproduced. A plan that is optimal when initially solved may no longer be optimal when future planners solve a plan of the same form from the forest states generated by earlier decisions. The central mechanism, as the thesis argued \citep[p.~134]{daugherty1991credibility}, is the interaction between the inter-period flow constraint and a replanning institution that rolls the horizon forward and re-anchors the constraint, treating realized harvests as sunk: when each replan instead solves the exact remainder of the original problem, the plan is followed. The phenomenon is not explained by non-unique LP optima or by a particular constant discount rate, and negatively valued strata contribute to it without being necessary for it. Of the alternative formulations tested, only replacing the flow constraint by a calibrated harvest cap nearly removes it.

The practical implication is that an open-loop trajectory should not be interpreted as a credible projection solely because it is optimal at the time the model is solved. Where planning involves sequential re-optimization, credibility requires that the announced plan remain optimal from the states it creates under the replanning institution actually used, or that the planning formulation explicitly account for subsequent replanning. The sequential-replanning framework developed here provides an operational means of testing this property and evaluating alternative formulations, as shown in extension experiments, against the trajectories they generate under repeated re-optimization. By providing an open and reproducible benchmark of Daugherty's analysis, this study makes dynamic consistency directly testable in forest-planning models and provides a basis for developing planning approaches whose projected outcomes remain credible under sequential implementation.